\documentclass[11pt,dvipdfmx,b5paper,oneside]{jsbook}

\usepackage{graphicx}
\usepackage{color}
\usepackage{here}
\usepackage{framed}
\usepackage{tcolorbox}
%\usepackage{quotchap}
\usepackage{pdfpages}
\usepackage[hidelinks]{hyperref}
\usepackage{pxjahyper}
\usepackage{titlesec}
\usepackage{picture}
\usepackage{tikz}

\tcbuselibrary{breakable}
\definecolor{shadecolor}{gray}{0.80}

% section
\titleformat{\section}[block]{}{}{0pt}
{
  \definecolor{teal}{gray}{0.30}
  \begin{picture}(0,0)
    \put(-10,-5){
      \begin{tikzpicture}
        \fill[teal] (0pt,0pt) rectangle (5pt,19pt);
      \end{tikzpicture}
    }
    \put(-10,-5){
      \color{teal}
      \line(1,0){\hsize}
    }
  \end{picture}
  \hspace{0pt}
  \sf \Large \thesection
  \hspace{0pt}
}

% 図表見出し
\renewcommand{\tablename}{\textcolor{gray}{▼} 表}
\renewcommand{\figurename}{\textcolor{gray}{▲} 図}

\begin{document}

\chapter[ニューラルネットワークのアーキテクチャの神経生理学的解釈]{ニューラルネットワークの\\アーキテクチャの神経生理学的解釈}
\section{畳み込みニューラルネットワーク (CNN)}
CNNは視覚のモデルとして優れていることは分かるのだが、スライディングウィンドウなどが入るせいでいまいち理解できない。
CNNはスライディングウィンドウを用いて視覚全体にニューロンをコピペしたもの。同種のニューロンは同じものとして重み共有により学習。
\subsection{max-pooling}
Winner-take allの仕組みはlateral inhibitionとして捉えることができる。
\section{LSTM}
セルアセンブリによるメモリの維持。他にはCA3の海馬の構造も似ている。

\section{GRU}

\section{Dropout}
Dropoutは負のノイズにより発火が抑制されること。Hintonのセミナー。\\
\url{https://www.youtube.com/watch?v=VIRCybGgHts}\\
\url{https://www.quora.com/Is-there-a-biological-basis-for-dropout-in-neural-networks}
正則化が如何に実現されているかは議論の余地があるが、ヒトのニューロンの特性が認知機能を支えているかもしれない。

\section{Normalization}
Normalization自体は昔からV1のモデルとして存在していた。
我輩は猫である\footnote{こんな感じで脚注を書く}。


\begin{tcolorbox}[breakable]
\begin{verbatim}
1  /* ここにはソースコードを書く */
2  #include<stdio.h>
3
4  int main(void)
5  {
6    printf("Hello, World!\n");
7    return 0;
8  }
9  /* breakableを付けるとこんな感じで改行にも対応できる */
\end{verbatim}
\end{tcolorbox}

\begin{shaded}
\begin{verbatim}
## ここにはコマンドを書く
$ echo "Hello, World!"
\end{verbatim}
\end{shaded}

図表はキャプションを付けたときに、先頭に「▲」や「▼」を付けるようにした。

\begin{table}[H]
  \centering
  \caption{表のサンプル}
  \begin{tabular}{|c|l|l|l|} \hline
    日本 & hoge & fuga & piyo \\ \hline
    アメリカ & foo & bar & baz \\ \hline
  \end{tabular}
  \label{table-sample}
\end{table}

\begin{tcolorbox}[title=抑制性介在ニューロンによる]
  コラムも随時挟めるようにした。

  tcolorboxはtitleを指定するといい感じにタイトル付きの枠で囲ってくれる。
\end{tcolorbox}
\begin{thebibliography}{99}
\subsection*{1章}
\bibitem{Woźniak} L. Beaulieu-Laroche, et al. ``Enhanced Dendritic Compartmentalization in Human Cortical Neurons.''. (2018). \url{https://www.ncbi.nlm.nih.gov/pubmed/30340039}
\end{thebibliography}

\end{document}